\subsection{California: Propositions 11 and 20}

\textbf{Proposition 11 (2008).} The California Redistricting Reform Act of 2008, authored primarily by California Common Cause, established the California Citizens Redistricting Commission (CCRC) with authority over state legislative districts. It passed with 50.9\% of the vote---the narrowest margin of the three case studies---in a year when the California Democratic Party spent over \$1 million against it.

The CCRC consists of 14 commissioners: five Democrats, five Republicans, and four members from neither major party. Commissioners are selected through an application and screening process administered by the State Auditor, with legislative leadership granted limited striking authority. No elected official, political consultant, or lobbyist may serve.

Proposition 11 required the CCRC to draw maps satisfying six criteria in priority order: (1) federal constitutional requirements, (2) Voting Rights Act compliance, (3) population equality, (4) geographic contiguity, (5) geographic compactness, and (6) respect for political subdivision boundaries. Incumbency protection and partisan advantage are not listed criteria and are prohibited considerations.

\textbf{Proposition 20 (2010).} Having established the CCRC for state legislative districts, reformers returned with Proposition 20, which extended CCRC authority to congressional districts. Proposition 20 passed with 61.1\% of the vote, substantially more comfortable than Proposition 11, in part because the CCRC's initial state legislative maps (drawn in 2011) had built public confidence. Proposition 27, on the same ballot, would have stripped the CCRC of congressional authority and returned it to the legislature; it failed with only 40.5\% support.

The combined effect of Propositions 11 and 20 is that California now has a single independent commission drawing both state and federal maps from the same criteria, using a single process, with a single audit mechanism. The 2021 redistricting cycle, the first under the combined authority, was extensively documented and received favorable assessments from good-government organizations for its public engagement and map quality.

\textbf{Lessons from California.} The narrow passage of Proposition 11 reflected low public awareness of redistricting; the comfortable passage of Proposition 20 two years later reflected both the legitimacy of the CCRC's early work and better-organized advocacy. The California experience suggests that initiatives can be self-reinforcing: a well-designed first initiative creates institutional legitimacy that makes subsequent reforms more popular. The failure of Proposition 27 in 2022 (after the CCRC's 2021 maps were widely regarded as fair) confirms this dynamic.

\subsection{Colorado: Amendments Y and Z}

Colorado's 2018 initiative campaign was notable for two features: it addressed congressional and legislative redistricting simultaneously, and it achieved overwhelming approval margins.

\textbf{Amendment Y} created an independent commission for congressional redistricting (12 members: 4 Democrats, 4 Republicans, 4 unaffiliated). \textbf{Amendment Z} created a parallel commission for state legislative redistricting (12 members with the same partisan structure). Both passed with approximately 71\% and 70\% support respectively---the highest margins of any redistricting reform initiative in the three-state sample.

Colorado's criteria structure differs from California's in one important respect: the commissions are explicitly prohibited from drawing maps with the intent to protect incumbents or to favor or disfavor any political party, and maps must not unduly favor any political party. This anti-favoritism criterion, absent in California's original Proposition 11, was added after observing that California's criteria-only approach left some ambiguity about intent-based versus effect-based partisan tests.

The Colorado commissions conducted the 2021 redistricting cycle with extensive public hearings (over 60 public meetings, including remote sessions during the COVID pandemic) and produced maps that were approved by the Colorado Supreme Court and implemented without federal litigation.

\textbf{Lessons from Colorado.} The simultaneous treatment of both chambers was logistically more complex but politically more defensible---it prevented the argument that legislative redistricting reform was undermined by the absence of congressional reform. The higher approval margins relative to California reflected both more effective public communication and Colorado's more competitive electoral environment (voters in competitive states perceive gerrymandering as costing them more directly). The explicit anti-favoritism criterion strengthens the legal foundation for challenging partisan manipulation.

\subsection{Michigan: Proposal 2}

Michigan's Proposal 2 (2018) created the Michigan Independent Citizens Redistricting Commission (MICRC) and was notable as the most aggressive structural reform of the three case studies: it amended the state constitution directly, placed the commission outside the control of the legislature or governor, and established strict conflict-of-interest rules that excluded not just elected officials but their family members.

The MICRC consists of 13 members: 4 Republicans, 4 Democrats, and 5 independents, selected through a process that begins with a random pool of applicants drawn from the Secretary of State's voter rolls (excluding ineligible persons). This selection-by-lottery mechanism---a form of sortition---is more neutral than California's application process and more resistant to partisan manipulation of the selection committee.

Michigan's criteria include explicit attention to communities of interest (groups of people who share common social or economic interests) and require the commission to consider the extent to which partisan fairness is achieved. The partisan fairness criterion was the subject of extensive litigation: in \textit{Agee v. Benson}, No.\ 21-2267 (6th Cir. 2022), challengers argued that the partisan fairness criterion violated the Elections Clause. The Sixth Circuit upheld the MICRC maps.

\textbf{Lessons from Michigan.} The sortition-based selection process is the most replicable innovation from the Michigan experience: it eliminates the bottleneck of politically controlled selection committees (the main vulnerability of California's State Auditor process, which was challenged in \textit{Rizo v. Padilla}). The broader community of interest definition explicitly acknowledges that redistricting affects groups beyond partisan blocs, which is important for VRA compliance. The partisan fairness criterion creates a tension with algorithmic neutrality that must be resolved in the model initiative language.
